% FAC-EN-U0046
% French source: fac_body.tex lines 1635--1680.
% Authority: PDF physical pp. 45--46; printed pp. 240--241 (section 4 opening and no. 48 complete).

\medskip
\phantomsection\label{fac:II:4}
\textbf{§4. Correspondence between finitely generated modules and coherent algebraic sheaves}

\phantomsection\label{fac:II:4:no:48}
\textbf{48. Sheaf associated with a module.}
Let $V$ be an affine variety and let $\mathcal{O}$ be the sheaf of local rings
\phantomsection\label{fac:II:4:no:48:def:coordinate-ring}
of $V$. The ring $A=\Gamma(V,\mathcal{O})$ will be called the
\textit{coordinate ring} of $V$; it is a $K$-algebra whose only nilpotent
element is $0$. If $V$ is embedded as a closed subvariety of an affine space
$K^r$, we know (cf. \hyperref[fac:II:3:no:44]{no.~44}) that $A$ is identified with the quotient of
$K[X_1,\ldots,X_r]$ by the ideal of polynomials vanishing on $V$. It follows
that the algebra $A$ is generated by finitely many elements.

Conversely, it is readily verified that if $A$ is a commutative $K$-algebra
without nilpotent elements other than $0$, and is generated by finitely many
elements, then there is an affine variety $V$ such that $A$ is isomorphic to
$\Gamma(V,\mathcal{O})$. Moreover, $V$ is determined up to isomorphism by
this property (one may identify $V$ with the set of characters of $A$,
endowed with the usual topology).

\phantomsection\label{fac:II:4:no:48:construction:A-M}
Let $M$ be an $A$-module. The module $M$ defines a constant sheaf on $V$,
which we shall again denote by $M$; similarly, $A$ defines a constant sheaf,
and the sheaf $M$ may be regarded as a sheaf of $A$-modules. Put
\[
\mathcal{A}(M)=\mathcal{O}\otimes_A M,
\]
where $\mathcal{O}$ is also regarded as a sheaf of $A$-modules. It is clear
that $\mathcal{A}(M)$ is an algebraic sheaf on $V$. Moreover, if
$\varphi:M\to M'$ is an $A$-homomorphism, there is a homomorphism
\[
\mathcal{A}(\varphi)=1\otimes\varphi:
\mathcal{A}(M)\longrightarrow\mathcal{A}(M').
\]
In other words, $M\mapsto\mathcal{A}(M)$ is a covariant functor.

\phantomsection\label{fac:II:4:prop:1}
\textsc{Proposition 1.} \textit{The functor
$M\mapsto\mathcal{A}(M)$ is exact.}

Let $M\to M'\to M''$ be an exact sequence of $A$-modules. We must show that
the sequence
$\mathcal{A}(M)\to\mathcal{A}(M')\to\mathcal{A}(M'')$ is exact; in other
words, that for every $x\in V$ the sequence
\[
\mathcal{O}_x\otimes_A M
\longrightarrow\mathcal{O}_x\otimes_A M'
\longrightarrow\mathcal{O}_x\otimes_A M''
\]
is exact.

Now $\mathcal{O}_x$ is precisely the ring of fractions $A_S$ of $A$, where
$S$ is the set of all $f\in A$ such that $f(x)\ne0$ (for the definition of a
ring of fractions, cf. [8], [12], or [13]). \hyperref[fac:II:4:prop:1]{Proposition~1} is therefore a
special case of the \hyperref[fac:II:4:lemma:1]{following result}.

\phantomsection\label{fac:II:4:lemma:1}
\textsc{Lemma 1.} \textit{Let $A$ be a ring, let $S$ be a multiplicatively
closed subset of $A$ not containing $0$, and let $A_S$ be the ring of
fractions of $A$ with respect to $S$. If $M\to M'\to M''$ is an exact
sequence of $A$-modules, then the sequence
\[
A_S\otimes_A M\longrightarrow A_S\otimes_A M'
\longrightarrow A_S\otimes_A M''
\]
is exact.}

Let $M_S$ denote the set of fractions $m/s$, with $m\in M$ and $s\in S$,
where two fractions $m/s$ and $m'/s'$ are identified if there is an
$s''\in S$ such that
\[
s''(s'\cdot m-s\cdot m')=0.
\]
It is readily seen that $M_S$ is an $A_S$-module and that the map
\[
a/s\otimes m\longmapsto a\cdot m/s
\]
is an isomorphism from $A_S\otimes_A M$ onto $M_S$. We are therefore reduced
to showing that the sequence
\[
M_S\longrightarrow M'_S\longrightarrow M''_S
\]
is exact, which is immediate.

\phantomsection\label{fac:II:4:prop:2}
\textsc{Proposition 2.} \textit{$\mathcal{A}(M)=0$ implies $M=0$.}

Let $m$ be an element of $M$. If $\mathcal{A}(M)=0$, then
$1\otimes m=0$ in $\mathcal{O}_x\otimes_A M$ for every $x\in V$. By the
preceding discussion, $1\otimes m=0$ is equivalent to the existence of an
  element $s\in A$, with $s(x)\ne0$, such that $s\cdot m=0$.
  The annihilator of $m$ is therefore contained in no maximal ideal of $A$; hence it is
equal to $A$, and consequently $m=0$.

\phantomsection\label{fac:II:4:prop:3}
\textsc{Proposition 3.} \textit{If $M$ is a finitely generated $A$-module,
$\mathcal{A}(M)$ is a coherent algebraic sheaf on $V$.}

Since $M$ is finitely generated and $A$ is Noetherian, $M$ is isomorphic to
the cokernel of a homomorphism $\varphi:A^q\to A^p$, and
$\mathcal{A}(M)$ is isomorphic to the cokernel of
\[
\mathcal{A}(\varphi):\mathcal{A}(A^q)\longrightarrow\mathcal{A}(A^p).
\]
Since $\mathcal{A}(A^p)=\mathcal{O}^p$ and
$\mathcal{A}(A^q)=\mathcal{O}^q$, it follows that $\mathcal{A}(M)$ is
coherent.
